\documentclass[10pt,a4paper]{article}

% Packages essentiels
\usepackage[utf8]{inputenc}
\usepackage[T1]{fontenc}
\usepackage[margin=1.5cm]{geometry}
\usepackage{amsmath,amssymb}
\usepackage{tikz}
\usetikzlibrary{positioning,calc,arrows.meta,shapes.geometric}
\usepackage{xcolor}
\usepackage{listings}
\usepackage{tcolorbox}
\tcbuselibrary{breakable}
\usepackage{hyperref}

% Configuration listings avec gestion complète des accents
\lstset{
    language=C,
    basicstyle=\ttfamily\scriptsize,
    backgroundcolor=\color{gray!10},
    keywordstyle=\color{blue},
    commentstyle=\color{gray},
    stringstyle=\color{red},
    numbers=left,
    numberstyle=\tiny,
    breaklines=true,
    frame=single,
    tabsize=4,
    showstringspaces=false,
    inputencoding=utf8,
    extendedchars=true,
    literate={é}{{\'e}}1 {è}{{\`e}}1 {à}{{\`a}}1 {ç}{{\c{c}}}1 
             {ê}{{\^e}}1 {î}{{\^i}}1 {ô}{{\^o}}1 {û}{{\^u}}1 {ù}{{\`u}}1
             {É}{{\'E}}1 {È}{{\`E}}1 {À}{{\`A}}1 {Ç}{{\c{C}}}1
             {Ê}{{\^E}}1 {Î}{{\^I}}1 {Ô}{{\^O}}1 {Û}{{\^U}}1 {Ù}{{\`U}}1
             {ë}{{\"e}}1 {ï}{{\"\i}}1 {ö}{{\"o}}1 {ü}{{\"u}}1
             {Ë}{{\"E}}1 {Ï}{{\"I}}1 {Ö}{{\"O}}1 {Ü}{{\"U}}1
             {œ}{{\oe}}1 {Œ}{{\OE}}1 {æ}{{\ae}}1 {Æ}{{\AE}}1
}

% Boxes pédagogiques
\newtcolorbox{exercicebox}[1]{
    colback=green!5,colframe=green!50!black,
    fonttitle=\bfseries,title=#1,breakable
}

\newtcolorbox{warningbox}[1]{
    colback=red!5,colframe=red!75!black,
    fonttitle=\bfseries,title=#1,breakable
}

\newtcolorbox{infobox}[1]{
    colback=blue!5,colframe=blue!75!black,
    fonttitle=\bfseries,title=#1,breakable
}

\title{\textbf{TD Chapitre 3 : Structures et Unions}\\
\large Programmation C Avancée - L3}
\author{Dr. El Hadji Bassirou TOURÉ\\
Département de Mathématiques et Informatique\\
Faculté des Sciences et Techniques\\
Université Cheikh Anta Diop de Dakar}
\date{}

\begin{document}

\maketitle

\begin{infobox}{Contexte}
Architecture x86-64, little-endian. Alignement: \texttt{char}=1, \texttt{short}=2, \texttt{int}=4, \texttt{long}/pointeur=8, \texttt{double}=8. Tous les dumps hexadécimaux utilisent 2 caractères par octet.
\end{infobox}

\section{Exercice 1 : Calcul de Padding}

Pour chaque structure, calculer: (a) l'offset de chaque membre, (b) le padding inséré, (c) \texttt{sizeof} total.

\begin{exercicebox}{1.1}
\begin{lstlisting}
struct A {
    char a;
    int b;
    char c;
};
\end{lstlisting}
\end{exercicebox}

\begin{exercicebox}{1.2}
\begin{lstlisting}
struct B {
    char a;
    char b;
    int c;
    short d;
};
\end{lstlisting}
\end{exercicebox}

\begin{exercicebox}{1.3}
\begin{lstlisting}
struct C {
    double x;
    char tag;
    int count;
    char flag;
};
\end{lstlisting}
\end{exercicebox}

\begin{exercicebox}{1.4}
\begin{lstlisting}
struct D {
    char a;
    short b;
    char c;
    long d;
    char e;
};
\end{lstlisting}
\end{exercicebox}

\section{Exercice 2 : Analyse de Dumps Hexadécimaux}

\begin{exercicebox}{2.1 - Reconstruction de structure}
Le dump suivant correspond à une variable de type \texttt{struct Data} à l'adresse \texttt{0x7FFE1000}:

\begin{verbatim}
0x7FFE1000: 41 00 00 00 78 56 34 12 42 43 00 00
\end{verbatim}

\begin{lstlisting}
struct Data {
    char tag;
    int value;
    short code;
};
\end{lstlisting}

Questions:
\begin{enumerate}
\item Identifier les octets de données et de padding
\item Reconstruire les valeurs: \texttt{tag}, \texttt{value}, \texttt{code}
\item Quel est le \texttt{sizeof(struct Data)} ?
\end{enumerate}
\end{exercicebox}

\begin{exercicebox}{2.2 - Structure imbriquée}
\begin{verbatim}
0x7FFE2000: 0A 00 00 00 14 00 00 00 1E 00 00 00 28 00 00 00
            52 00 00 00 00 00 00 00
\end{verbatim}

\begin{lstlisting}
struct Point { int x, y; };
struct Rectangle {
    struct Point top_left;
    struct Point bottom_right;
    char color;
};
\end{lstlisting}

Questions:
\begin{enumerate}
\item Identifier chaque champ dans le dump
\item Extraire les coordonnées des deux points
\item Quelle est la valeur de \texttt{color} (code ASCII) ?
\item Combien d'octets de padding final ?
\end{enumerate}
\end{exercicebox}

\begin{exercicebox}{2.3 - Tableau de structures}
\begin{verbatim}
0x7FFE3000: 41 00 64 00 42 00 C8 00 43 00 2C 01
\end{verbatim}

\begin{lstlisting}
struct Entry {
    char id;
    short value;
};
struct Entry entries[3];
\end{lstlisting}

Questions:
\begin{enumerate}
\item Quel est le \texttt{sizeof(struct Entry)} ?
\item Reconstruire le tableau \texttt{entries}
\item Quel est l'offset de \texttt{entries[2].value} ?
\end{enumerate}
\end{exercicebox}

\section{Exercice 3 : Optimisation du Padding}

\begin{exercicebox}{3.1 - Réorganisation}
\begin{lstlisting}
struct Original {
    char a;
    double b;
    char c;
    int d;
    char e;
};
\end{lstlisting}

Questions:
\begin{enumerate}
\item Calculer \texttt{sizeof(struct Original)}
\item Proposer une réorganisation minimisant la taille
\item Calculer la nouvelle taille
\item Économie réalisée en octets et en pourcentage
\end{enumerate}
\end{exercicebox}

\begin{exercicebox}{3.2 - Structure avec pointeurs}
\begin{lstlisting}
struct Node {
    char type;
    struct Node *next;
    int data;
    char flags;
};
\end{lstlisting}

Questions:
\begin{enumerate}
\item Calculer \texttt{sizeof(struct Node)} (x86-64)
\item Proposer une version optimisée
\item Quelle taille pour la version optimisée ?
\end{enumerate}
\end{exercicebox}

\section{Exercice 4 : Unions}

\begin{exercicebox}{4.1 - Interprétation mémoire}
\begin{lstlisting}
union U {
    int i;
    float f;
    char c[4];
};

union U u;
u.i = 0x40490FDB;
\end{lstlisting}

Questions:
\begin{enumerate}
\item Quel est le dump hexadécimal de \texttt{u} (little-endian) ?
\item Quelle est la valeur de \texttt{u.c[0]}, \texttt{u.c[1]}, \texttt{u.c[2]}, \texttt{u.c[3]} ?
\item \texttt{0x40490FDB} est la représentation IEEE 754 de $\pi$. Quelle sera la valeur (approximative) de \texttt{u.f} ?
\end{enumerate}
\end{exercicebox}

\begin{exercicebox}{4.2 - Analyse de dump}
\begin{verbatim}
0x7FFE4000: 01 00 00 00 00 00 F0 3F
\end{verbatim}

\begin{lstlisting}
typedef struct {
    int type;
    union {
        int i;
        double d;
    } value;
} Variant;
\end{lstlisting}

Questions:
\begin{enumerate}
\item Quel est \texttt{sizeof(Variant)} ?
\item Quelle est la valeur de \texttt{type} ?
\item Si \texttt{type == 1} signifie \texttt{double}, quelle est la valeur stockée ?

Note: \texttt{0x3FF0000000000000} est la représentation IEEE 754 de \texttt{1.0}.
\end{enumerate}
\end{exercicebox}

\begin{exercicebox}{4.3 - Détection d'endianness}
\begin{lstlisting}
union EndianCheck {
    int i;
    char c[4];
};

int is_little_endian(void) {
    union EndianCheck e;
    e.i = 1;
    return e.c[0] == 1;
}
\end{lstlisting}

Questions:
\begin{enumerate}
\item Expliquer pourquoi ce code fonctionne
\item Dessiner le layout mémoire de \texttt{e} pour little-endian
\item Dessiner le layout mémoire de \texttt{e} pour big-endian
\end{enumerate}
\end{exercicebox}

\section{Exercice 5 : Bit Fields}

\begin{exercicebox}{5.1 - Calcul de taille}
\begin{lstlisting}
struct Flags1 {
    unsigned int a : 1;
    unsigned int b : 3;
    unsigned int c : 4;
};

struct Flags2 {
    unsigned int a : 1;
    unsigned int b : 3;
    unsigned int : 0;  // force alignment
    unsigned int c : 4;
};
\end{lstlisting}

Questions:
\begin{enumerate}
\item Quel est \texttt{sizeof(struct Flags1)} ?
\item Quel est \texttt{sizeof(struct Flags2)} ?
\item Expliquer l'effet de \texttt{: 0}
\end{enumerate}
\end{exercicebox}

\begin{exercicebox}{5.2 - Représentation binaire}
\begin{lstlisting}
struct Status {
    unsigned int active : 1;
    unsigned int mode : 2;
    unsigned int priority : 3;
    unsigned int reserved : 2;
};

struct Status s = {1, 2, 5, 0};
\end{lstlisting}

Questions:
\begin{enumerate}
\item Représenter \texttt{s} en binaire (8 bits)
\item Quelle est la valeur hexadécimale du premier octet ?
\item Si on lit \texttt{s} comme un \texttt{unsigned char}, quelle valeur obtient-on ?
\end{enumerate}
\end{exercicebox}

\section{Exercice 6 : Structures et Allocation}

\begin{exercicebox}{6.1 - Analyse de code}
\begin{lstlisting}
struct Person {
    char *name;
    int age;
};

struct Person *p = malloc(sizeof(struct Person));
p->name = "Alice";
p->age = 25;
// ...
free(p);
\end{lstlisting}

Questions:
\begin{enumerate}
\item Ce code est-il correct ? Justifier.
\item Proposer une version avec allocation de \texttt{name}
\item Écrire la fonction de libération correspondante
\end{enumerate}
\end{exercicebox}

\begin{exercicebox}{6.2 - Flexible Array Member}
\begin{lstlisting}
typedef struct {
    size_t len;
    char data[];
} String;

String *s = malloc(sizeof(String) + 6);
s->len = 5;
strcpy(s->data, "Hello");
\end{lstlisting}

Questions:
\begin{enumerate}
\item Quel est \texttt{sizeof(String)} ?
\item Pourquoi alloue-t-on \texttt{sizeof(String) + 6} ?
\item Dessiner le layout mémoire de \texttt{s}
\end{enumerate}
\end{exercicebox}

\begin{exercicebox}{6.3 - Structure opaque}
Soit le header:
\begin{lstlisting}
// stack.h
typedef struct Stack Stack;
Stack *stack_create(size_t capacity);
void stack_push(Stack *s, int value);
int stack_pop(Stack *s);
void stack_destroy(Stack *s);
\end{lstlisting}

Questions:
\begin{enumerate}
\item Pourquoi ne peut-on pas écrire \texttt{Stack s;} dans le code client ?
\item Pourquoi peut-on écrire \texttt{Stack *s;} ?
\item Proposer une implémentation de \texttt{struct Stack}
\end{enumerate}
\end{exercicebox}

\section{Exercice 7 : Forensics Avancé}

\begin{exercicebox}{7.1 - Structure complexe}
\begin{verbatim}
0x7FFE5000: 01 00 00 00 00 00 00 00 00 10 FE 7F 00 00 00 00
            05 00 00 00 00 00 00 00
\end{verbatim}

\begin{lstlisting}
typedef struct {
    int type;
    void *data;
    size_t size;
} Buffer;
\end{lstlisting}

Questions:
\begin{enumerate}
\item Calculer \texttt{sizeof(Buffer)} et identifier le padding
\item Extraire les valeurs de \texttt{type}, \texttt{data}, \texttt{size}
\item À quelle adresse pointe \texttt{data} ?
\end{enumerate}
\end{exercicebox}

\begin{exercicebox}{7.2 - Tableau dynamique de structures}
\begin{verbatim}
0x00602000: 00 00 00 00 0A 00 00 00 01 00 00 00 14 00 00 00
            02 00 00 00 1E 00 00 00
\end{verbatim}

\begin{lstlisting}
struct Item {
    int id;
    int value;
};

struct Item *items = malloc(3 * sizeof(struct Item));
\end{lstlisting}

Questions:
\begin{enumerate}
\item Reconstruire le contenu du tableau
\item Quel est l'offset de \texttt{items[1].value} par rapport à \texttt{items} ?
\item Écrire l'expression pour accéder à \texttt{items[2].id} via arithmétique de pointeurs
\end{enumerate}
\end{exercicebox}

\end{document}
